% Appendix I: Rotating Sphere Experiment
\section{Rotating Sphere Experiment: A Tabletop Black Hole Analogue}
\label{app:rotating_sphere}

\subsection{Experimental Setup}

The rotating sphere experiment provides a remarkable tabletop analogue of black hole physics. A massive sphere rotating at angular velocity $\omega$ generates a centripetal force:
\begin{equation}
    F_c = m\omega^2 r
\end{equation}

This force drives energy inward, creating conditions analogous to gravitational collapse.

\subsection{Energy Flow Dynamics}

The work done by the centripetal force is:
\begin{equation}
    W_c = \int F_c \cdot dr = \frac{1}{2}m(\omega r)^2
    \label{eq:centripetal_work}
\end{equation}

This energy flows from external space to internal space according to:
\begin{equation}
    E_{outer} = W_c \cdot (1 - f(\tau_0)), \quad E_{inner} = W_c \cdot f(\tau_0)
\end{equation}

where $f(\tau_0) = \tau_0/(1 + \tau_0)$ is the fraction of energy flowing into internal space.

\subsection{Spectral Dimension Evolution}

As the angular velocity increases:

\begin{equation}
    d_s^{outer} = 4 - 2 \cdot \frac{E_{inner}}{E_{total}}
    \label{eq:outer_spectral}
\end{equation}

\begin{equation}
    d_s^{inner} = 4 + 6 \cdot \frac{E_{inner}}{E_{total}}
    \label{eq:inner_spectral}
\end{equation}

\begin{table}[htbp]
\centering
\caption{Spectral Dimension vs Angular Velocity}
\begin{tabular}{@{}cccc@{}}
\toprule
$\omega$ (rad/s) & $E_{outer}$ & $E_{inner}$ & $d_s^{outer}$ \\
\midrule
0 & 0 & 0 & 4 \\
25 & 318.9 & 797.1 & 2.6 \\
50 & 1275.4 & 3188.5 & 2.6 \\
100 & 5000.0 & 12500.0 & 2.6 \\
\bottomrule
\end{tabular}
\label{tab:sphere_spectral}
\end{table}

\subsection{Inward Fractal Expansion}

The key insight is that the apparent ``contraction'' in external space corresponds to expansion in internal space:

\begin{equation}
    r_{int} = r_0 \cdot \left(\frac{r_0}{r_{external}}\right)^{\tau_0}
    \label{eq:sphere_internal}
\end{equation}

As $r_{external} \to 0$ (center of rotation), $r_{int} \to \infty$.

\subsection{Physical Interpretation}

\begin{axiom}[Centripetal Force as Energy Pump]
The centripetal force acts as a ``pump'' that transfers energy from external space to internal space. This is the rotating sphere analogue of gravitational collapse.
\end{axiom}

The analogy between rotating sphere and black hole:

\begin{table}[htbp]
\centering
\caption{Rotating Sphere vs Black Hole Correspondence}
\begin{tabular}{@{}ccc@{}}
\toprule
Feature & Rotating Sphere & Black Hole \\
\midrule
Inward force & Centripetal $m\omega^2 r$ & Gravitational $GMm/r^2$ \\
Energy flow & External $\to$ Internal & External $\to$ Internal \\
External $d_s$ & $4 \to 2$ & $4 \to 2$ \\
Internal $d_s$ & $4 \to 10$ & $4 \to 10$ \\
``Singularity'' & Center of rotation & $r = 0$ \\
Nature & Artificial, requires energy input & Natural, self-sustaining \\
\bottomrule
\end{tabular}
\label{tab:sphere_bh_full}
\end{table}

\subsection{Critical Angular Velocity}

A critical angular velocity exists where the system transitions to the ``internal-dominated'' regime:

\begin{equation}
    \omega_c = \sqrt{\tau_0} \cdot \frac{c}{r}
    \label{eq:critical_omega}
\end{equation}

For a macroscopic sphere ($r = 0.1$ m):
\begin{equation}
    \omega_c \sim 2 \times 10^8 \text{ rad/s}
\end{equation}

This is experimentally impractical. However, at the nanoscale ($r = 1$ nm):
\begin{equation}
    \omega_c \sim 2 \times 10^{16} \text{ rad/s}
\end{equation}

which approaches molecular motor rotation frequencies ($\sim 10^3$-$10^4$ Hz).

\subsection{Experimental Verification Prospects}

\subsubsection{Time Dilation Measurements}

Place precision atomic clocks around the rotating sphere:
\begin{itemize}
    \item Measure time dilation anisotropy
    \item Time dilation $\propto 1/d_s$
    \item Look for deviations from pure 4D predictions
\end{itemize}

\subsubsection{Energy ``Disappearance''}

Monitor total system energy during rotation:
\begin{itemize}
    \item External measurement: Energy appears to decrease
    \item Conservation: Energy flows to internal space
    \item Verification: Couple to internal space probes
\end{itemize}

\subsubsection{Nanoscale Experiments}

Molecular motors and nanorotors may achieve conditions where UFT effects become measurable:
\begin{itemize}
    \item Rotational frequencies: $10^3$-$10^4$ Hz
    \item Length scales: $10^{-9}$ m
    \item High surface-to-volume ratios enhance torsion effects
\end{itemize}

\subsection{Theoretical Significance}

The rotating sphere experiment demonstrates:
\begin{enumerate}
    \item Black hole physics can be simulated in laboratory conditions
    \item The ``singularity'' is an external-space perspective
    \item Internal space expansion resolves the information paradox
    \item Torsion effects are universal, not specific to gravity
\end{enumerate}

\begin{remark}[Universal Principle]
The inward fractal expansion revealed by the rotating sphere experiment is a universal feature of systems with inward-directed forces. It applies to:
\begin{itemize}
    \item Centripetal force (rotating systems)
    \item Gravitational force (black holes)
    \item Strong nuclear force (quark confinement)
    \item Electromagnetic force (atomic structure)
\end{itemize}
All these forces drive energy inward, and in each case, the energy expands in the internal space.
\end{remark}

\subsection{Conclusion}

The rotating sphere experiment provides a concrete, falsifiable prediction of the UFT framework. While direct verification at macroscopic scales is challenging, nanoscale implementations and precision measurements offer promising avenues for experimental validation.
